% Lösungen Einheit 2 von 4 – Proportionale Zuordnungen und Dreisatz (zusammenbau v0.1)
\einheitenkopf{Einheit 2 von 4 · Proportionale Zuordnungen und Dreisatz}

\erg{12}{a) ein Kilogramm Äpfel}
\erg{13}{a) links: 5; 1; 8 (kg)}
\erg{14}{a) rechts ebenfalls : 8}
%% TODO Lösung der Erklärzeile 15g) fehlt in der Bank
\erg{15}{a) 6 €; 8 € \quad b) 6 € \quad c) 1 € \quad d) 6 km \quad e) 2 € \quad f) 30 m² \quad g) \ldots \quad h) 9 € \quad i) 15 € \quad j) 6 : 3 = 2 €; 2 · 7 = 14 € \quad k) 3,50 : 5 = 0,70 €; 0,70 · 9 = 6,30 € \quad l) Packung B: 100 g kosten 0,36 € statt 0,40 €. \quad m) 13 : 2 = 6,5 min je km; 6,5 · 5 = 32,5 min}
\erg{16}{a) 2,50 € je kg; y = 2,50 · x}
\erg{17}{a) 7,50 €; Gerade durch (0|0) und (1|3)}
\erg{18}{a) $\frac{x}{10} = \frac{300}{4}$, also x = 750 g}
\erg{19}{a) Tim hat subtrahiert statt geteilt: 7,20 : 3 = 2,40 € je kg. \quad b) Bei proportionalen Zuordnungen gehört zu einem Viertel der Menge ein Viertel des Preises – beide Seiten ändern sich im gleichen Verhältnis. \quad c) 1,50 € je kg; y = 1,50 · x \quad d) 250 : 4 · 10 = 625 g; nein, eine Packung mit 500 g reicht nicht.}
